\documentclass[a4paper, 11pt]{article}
    \usepackage[utf8]{inputenc}
    \usepackage[czech]{babel}
    \usepackage[text={17.2cm, 24.7cm}, left=2cm, top=2cm]{geometry}
    \usepackage{times}
    \usepackage[hidelinks]{hyperref}
    \usepackage{indentfirst}
    \usepackage{graphicx}
        \graphicspath{{./images/}}
    \usepackage{float}
    \usepackage{xcolor}
    \usepackage{enumitem}

\newcommand{\image}[3]{ % args: size [factor of text width], path to image, caption
    \begin{figure}[H]
        \centering
        \includegraphics[width=#1\textwidth]{#2}\caption{#3}
        \label{fig:\thefigure}
    \end{figure}
}

\renewcommand{\t}[1]{\texttt{#1}} % shortcut for tt text

\newcommand{\sss}[1]{\subsubsection{#1}} % shortcut for s.s.sec.

    \title {
        {\Huge \emph{\textbf{resipi.}}{\huge}} \\[0.25em] %, Kuchařka nové generace
        {\Large \emph{Zadání a návrh aplikace pro projekt do předmětu ITU}}
    }
    \author {
        \textbf{{\Large Tomáš Bordák,}
        \emph{\large xborda01}} \\[0.5em]
        {\Large Martin Jabůrek,} 
        \emph{\large xjabur02} \\[0.3em]
        {\Large Ondřej Šatinský,} 
        \emph{\large xsatin03} \\[0.5em]
        {\Large Tomáš Tomcsányi,} 
        \emph{\large xtomcs00}
    }
    \date{\today}

\begin{document}

    \maketitle

    \tableofcontents
    \listoffigures
    \newpage


    \section{Téma aplikace a motivace}

    Představíme-li si kuchařku, možná se nám vybaví babiččin volně svázaný sešit se všemi jejími
    osvědčenými recepty. S dnešní rozšířeností a dostupností počítačů, jako jsou i chytré telefony,
    se ale takové řešení už nejeví příliš atraktivní. Téměř každý už nosí v kapse přístroj schopný
    ukládat a zpracovávat snad jakákoli data i bez nepořádku a nepřehlednosti spojené s typickou
    papírovou variantou.

    Vařit může každý. Zároveň ale nikdo nevaří sám. Lidé si mezi sebou věky předávali své znalosti.
    Naučené recepty podle svého vkusu či umu upravovali a vylepšovali. Kdyby jsme si své recepty
    mezi sebou nesdíleli, kdo ví kde by kuchařské umění v dnešní době bylo. Opět se tak dostáváme
    k formě uchování receptů. Služby od emailu až po celosvětové sociální sítě nám umožňují
    komunikovat rychle a snadno jako nidky předtím. Už nemáme důvod vracet se k přepisování
    receptů na kus papíru nebo diktování jich do sluchátka telefonu.

    Mluvíme-li o sociálních sítích, jejich záběr je překonán jen málo čím. Miliony lidí je denně
    používájí k vyhledávání nových a často i překvapivých informací. A to i receptů.

    Cílem našeho týmu je zhotovit aplikaci, která nabídne nový pohled na uchování,
    sdílení a vaření podle receptů s využitím moderních informačních technologí.
    Tuto funkcionalitu spojí do moderní a chytré kuchařky,
    ve které si kterýkoli, ať už zkušený, či teprve začínající kuchař, najde
    přesně to, co jej zajímá.


    \section{Zadání}


    \subsection{Uživatelský průzkum a specifikace}
    
    \sss{Kdo je konkrétní uživatel}

    Aplikace je mířena zejména na studenty středních a vysokých škol. Právě tato věková kategorie
    nejen vaří, ale hlavně má nejvíce zkušeností s různými sociálními sítěmi. Je však
    nutné zmínit, že aplikace se na tuto skupinu neomezuje. V uživatelském průzkumu
    byl kladen důraz na co nejširší záběr. Je tedy kladen ohled i na zkušenější kuchaře, typicky
    vyššího věku, co se ale nemusí tolik orientovat v informačních technologiích.
    Aplikace by měla být dobře využitelná každým, kdo chce vařit.

    \sss{Co přesně uživatel od aplikace potřebuje}

    Uživatelé potřebují snadným a jednotným způsobem uchovávat, hledat a vyhledávat mezi recepty. Kromě tohoto
    potřebují recepty sdílet se svými blízkými. Pro usnadnění průběhu vaření dále uživatelé
    potřebují interaktivní a intuitivní způsob, jak se při vaření řídit receptem.

    Každý člen týmu provedl vlastní průzkum. V podkapitolách níže jsou popsány jednotlivé metody
    a závěry jimi zjištěné. Přesné potřeby uživatelů byly vyhodnoceny na základě tohoto průzkumu.

    \sss{Požadavky uživatele}

    % vychází z průzkumu
    \sss{Uživatelský průzkum Tomáše B.}
    \sss{Uživatelský průzkum Martina}
    \sss{Uživatelský průzkum Ondřeje}
    \sss{Uživatelský průzkum Tomáše T.}
    % shrnutí průzkumu každého každého člena
    \sss{Celkové shrnutí uživatelského průzkumu\,--\,výsledná specifikace}
    % celkové shrnutí

    \subsection{Průzkum existujících řešení}

    \sss{Existujícící aplikace řešicí požadovanou funkci}

    % celkový seznam

    \sss{Průzkum aplikací Tomáše B.}
    \sss{Průzkum aplikací Martina}
    \sss{Průzkum aplikací Ondřeje}
    \sss{Průzkum aplikací Tomáše T.}

    % individuální průzkumy, 2 aplikace, 3x +, 3x - (vč. inspirací, absencí)

    \sss{Závěr průzkumu}

    % celkový závěr, 3 inspirace/ponaučení


    \subsection{Zadání}

    % popis cílové app.
    % klíčové vlastnosti FE

    \subsection{Rozdělení práce na FE}

    % viz. dc
    % (máme stejný framework, každý má jistou rozsáhlou část stejné app.)
    % Úkoly
    %% každého člena
    %% týmové
    % u nás jednotlivé prvky aplikace


    \section{Návrh}


    \subsection{Návrh GUI}

    \sss{Návrh Tomáše B.}
    \sss{Návrh Martina}
    \sss{Návrh Ondřeje}
    \sss{Návrh Tomáše T.}

    % popisy návrhů
    % jak jsou řešeny požadavky uživatele
    % včetně obrázků


    \subsection{Výběr technologií}

    % FE (vue)
    % popř. BE, služby (...)
    % ZDŮVODNĚNÍ


    \subsection{Náveh API k BE} 

    % popis API (klíčové funkce !), popř. knihovny
    % klíčové datové struktury modelu
    % napojení BE na FE (tzn API na GUI, které funkce dělají co)

    %%%








    %\sss{Osobní interview}

    Každý člen týmu vybral skupinu potenciálních uživatelů, se kterou individuálně provedl interview.
    Účelem tohoto šetření bylo určit potřeby každé kategorie osob, které je pak nutné reflektovat
    ve výsledné aplikaci. Rozhovory byly provedeny s lidmi různých věkových kategorií. Jednalo se
    zejména o naše přátele a sourozence, ale i rodiče, kteří zastupovali starší uživatele.
    Zároveň byli uživatelé zvoleni tak, aby měli různou zkušenost s používáním informačních
    technologií. Jednalo se o uživatele s málo zkušnenostmi (typicky rodiče), každodenní uživatele
    mobilních zařízení (kamarádi a sourozenci) a dokonce i uživatele s oborovými znalostmi
    v informatice.

    Pro interview byla sestavena sada otázek, které se každý člen týmu držel:

    \begin{enumerate}[itemsep=0.25em, parsep=0em, label=\arabic*)]
        \item Jak často vaříš?
        \item Používáš při vaření kuchařku?
        \item Hledáš recepty na internetu?
        \item Podle čeho vyhledáváš?
        \item Co se ti na tom líbí / nelíbí?
        \item Kdyby existovala nějaká aplikace, která by ti umožnila tyhle věci vyhledávat, používal/a bys ji?
        \item Co by si v ní chtěl/a dělat?
        \item Proč?
        \item Máš přátele se kterými sdílíš recepty?
        \item Uvítal/a bys, kdyby to taková aplikace zjednodušila?
        \item Jaké používáš sociální sítě?
        \item Používáš tuto síť k vyhledávání receptů?
    \end{enumerate}

    % shrnutí získaných znalostí
    % ? uvedení výsledků výzkumu v příloze ?

    % Specifikace vychází z uživateského průzkumu

    %\sss{Osobní zhodnocení existujících řešení}

    Výše popsanému interview vždy následovala \uv{praktická} část, kde jsme potenciálním uživatelům
    předložili sadu podobných aplikací zvolených v \ref{pruz_exist_res}. Každý z nás dal svým
    uživatelům na vyzkoušení vlastní zvolené aplikace.

    % opět shrnutí získaných znalostí, individuální věc
    % ? potenciální příloha s kompletním průzkumem ?

    %\sss{Online dotazník}
    Aplikace je mířena zejména na studenty středních a vysokých škol. Právě tato věková kategorie
    nejen vaří, ale hlavně má nejvíce zkušeností s různými sociálními sítěmi. Je však
    nutné zmínit, že aplikace se na tuto skupinu neomezuje. V uživatelském průzkumu
    byl kladen důraz na co nejširší záběr. Je tedy kladen ohled i na zkušenější kuchaře, typicky
    vyššího věku, co se ale nemusí tolik orientovat v informačních technologiích.
    Aplikace by měla být dobře využitelná každým, kdo chce vařit.
    %\subsection{Průzkum existujících řešení}
    \label{pruz_exist_res}

    Každý člen týmu vyhledal sadu aplikací řešící podobný problém 
    % sada existicujících aplikací splňujících zadadou funkci
    %% 3 konkrétní +, -, popř. chybějící prvky každé
    % závěr
    %% za každého člena
    %% alespoň 3 inspirace, ponaučení

    %%%

\end{document}
